%==============================================================================
% FICHAMENTO CRÍTICO — Enhancing Deep Learning Classification Performance of
% Tongue Lesions in Imbalanced Data: Mosaic-Based Soft Labeling with
% Curriculum Learning
% Conforme NBR 6023 (referências) e NBR 10520 (citações)
%==============================================================================
\section{Fichamento: \citeonline{lee2024tongue}}

% --- 1. IDENTIFICAÇÃO ---
\subsection{Identificação}
\begin{tabular}{ll}
\textbf{Título:} & Enhancing deep learning classification performance of tongue
lesions in imbalanced data: mosaic-based soft labeling with curriculum learning \\
\textbf{Autor(es):} & Lee, Sung-Jae; Oh, Hyun Jun; Son, Young-Don; Kim, Jong-Hoon;
Kwon, Ik-Jae; Kim, Bongju \\
\textbf{Periódico:} & BMC Oral Health \\
\textbf{Ano:} & 2024 \\
\textbf{Volume:} & 24 \\
\textbf{Número:} & 1 \\
\textbf{Páginas:} & 161 \\
\textbf{DOI:} & \url{https://doi.org/10.1186/s12903-024-03898-3} \\
\textbf{Qualis:} & A2 (Odontologia) \\
\textbf{Tipo:} & Artigo original \\
\end{tabular}

% --- 2. OBJETIVO ---
\subsection{Objetivo}
Avaliar uma nova técnica de inteligência artificial baseada em mosaic-based soft
labeling com curriculum learning (CL) para melhorar o desempenho de classificação
de lesões de língua em um conjunto de dados desbalanceado, utilizando o modelo
EfficientNet-B0 como extrator de características e Grad-CAM para visualização das
regiões relevantes.

% --- 3. METODOLOGIA ---
\subsection{Metodologia}
\textbf{Desenho do estudo:} Estudo experimental de desenvolvimento e validação de
modelo de deep learning. \\
\textbf{Amostra:} 1.810 imagens de língua, sendo 372 de câncer, 141 de OPMDs
(desordens potencialmente malignas orais) e 1.297 de lesões não cancerosas. \\
\textbf{Métricas principais:} Acurácia, precisão, recall, F1-score, análise por
curva ROC. \\
\textbf{Validação:} Divisão treino/validação/teste (1.630 treino + 1.200 imagens
sintéticas / 180 teste), com validação cruzada interna e comparação com métodos
baseline (oversampling e weight balancing).

% --- 4. RESULTADOS PRINCIPAIS ---
\subsection{Resultados Principais}
\begin{itemize}
  \item O modelo final (MBO: mosaic-based soft labeling + curriculum learning)
        alcançou acurácia de 0,9444, superando oversampling (0,9111) e weight
        balancing (0,9278).
  \item A precisão relativa para a classe minoritária OPMD melhorou 21,2\% (de
        0,6875 para 0,8333) e o F1-score relativo melhorou 4,9\% (de 0,7333 para
        0,7692).
  \item O Grad-CAM foi utilizado para extrair regiões representativas de cada
        imagem, convertidas em mapas binários via limiar de Otsu para compor as
        imagens mosaic sintéticas.
\end{itemize}

% --- 5. LIMITAÇÕES ---
\subsection{Limitações}
\begin{itemize}
  \item O desempenho global ficou aquém das expectativas iniciais dos autores,
        com melhora concentrada apenas na classe minoritária OPMD.
  \item A amostra total é modesta (1.810 imagens) e proveniente de um único
        centro, limitando a validade externa dos resultados.
  \item O método proposto aumenta a complexidade computacional do pipeline de
        treinamento (duas etapas com geração de imagens sintéticas intermediárias).
\end{itemize}

% --- 6. CONTRIBUIÇÃO PARA O LIVRO ---
\subsection{Contribuição para o Livro}
Este artigo fundamenta diretamente os Capítulos 7 e 10, que tratam do uso de
EfficientNet com Grad-CAM para classificação de OPMDs. A estratégia de
mosaic-based soft labeling com CL oferece uma solução concreta para o problema
de desbalanceamento de classes, recorrente em bases de dados odontológicas
(onde lesões malignas e OPMDs são muito menos frequentes que lesões benignas).
O uso do Grad-CAM como mecanismo de explainable AI (XAI) para guiar a
augmentação de dados é uma contribuição metodológica relevante para o arcabouço
SDD/TDD do livro. O código e a abordagem são replicáveis e podem ser integrados
ao pipeline de experimentação descrito na Parte 3 do livro.

% --- 7. ANÁLISE CRÍTICA ---
\subsection{Análise Crítica}
\begin{itemize}
  \item \textbf{Validade interna:} Alta -- O estudo controla adequadamente as
        variáveis de confusão (arquitetura fixa, transfer learning com ImageNet,
        comparação com baselines estabelecidos).
  \item \textbf{Validade externa:} Média -- Amostra de um único centro sul-coreano,
        sem validação em populações multiétnicas ou multicêntricas.
  \item \textbf{Reprodutibilidade:} Alta -- A metodologia é descrita com detalhes
        suficientes (EfficientNet-B0, Grad-CAM, Otsu, CL em duas etapas) para
        replicação independente.
  \item \textbf{Relevância clínica:} Alta -- OPMDs são lesões de difícil
        classificação mesmo para especialistas; a melhoria de 21,2\% na precisão
        pode ter impacto direto no encaminhamento e biópsia.
  \item \textbf{Conflito de interesses:} Não -- Os autores declararam ausência de
        conflitos.
  \item \textbf{Viés identificado:} Viés de seleção -- a amostra é oriunda de
        registros hospitalares, podendo não representar a distribuição real de
        lesões de língua na população geral.
\end{itemize}

% --- 8. CITAÇÕES RELEVANTES ---
\subsection{Citações Relevantes}
\begin{citacao}
``The final model (MBO), which was trained using the newly added mosaic-based
soft labeling dataset (Stage 2), achieved an accuracy rate of 0.9444. This
outperformed conventional oversampling (0.9111) and weight balancing methods
(0.9278).'' (p. 11).
\end{citacao}

% --- 9. MAPEAMENTO SDD ---
\subsection{Mapeamento SDD}
\textbf{Problema:} Classificação de lesões de língua (câncer, OPMD, não cancerosa)
em dados desbalanceados com baixa representatividade da classe OPMD. \\
\textbf{Entrada:} Imagens fotográficas de língua (RGB), 1.810 amostras com
desequilíbrio de classes (372 câncer, 141 OPMD, 1.297 não cancerosas). \\
\textbf{Saída:} Probabilidades por classe (câncer, OPMD, não cancerosa) com
mapeamento por mosaic-based soft labeling (rótulo contínuo por área ocupada). \\
\textbf{Critérios:} Acurácia $\geq$ 0,94; precisão OPMD $\geq$ 0,83; F1-score
OPMD $\geq$ 0,76; superação dos baselines oversampling e weight balancing. \\
\textbf{Confiança:} 0,75 -- Estudo bem delineado, mas com amostra unicêntrica
modesta e resultados aquém das expectativas iniciais dos autores.

% --- 10. REFERÊNCIA ABNT ---
\subsection{Referência ABNT}
\begin{verbatim}
LEE, Sung-Jae et al. Enhancing deep learning classification performance of
tongue lesions in imbalanced data: mosaic-based soft labeling with curriculum
learning. BMC Oral Health, London, v. 24, n. 1, p. 161, 2024. DOI:
10.1186/s12903-024-03898-3.
\end{verbatim}
